\documentclass[11pt]{article}
\usepackage[margin=1.1in]{geometry}
\usepackage{amsmath,amssymb,amsthm}



\usepackage{enumitem}
\usepackage[colorlinks=true,linkcolor=blue,citecolor=blue,urlcolor=blue]{hyperref}

\newtheorem{theorem}{Theorem}[section]
\newtheorem{proposition}[theorem]{Proposition}
\newtheorem{conjecture}[theorem]{Conjecture}
\newtheorem{remark}[theorem]{Remark}
\theoremstyle{definition}
\newtheorem{definition}[theorem]{Definition}

\newcommand{\pP}{\textbf{[P]}}
\newcommand{\pC}{\textbf{[C]}}
\newcommand{\pO}{\textbf{[O]}}
\newcommand{\Tr}{\operatorname{Tr}}

\title{\textbf{Entanglement as the Origin of Spacetime Curvature:\\
A Verified Synthesis of the First-Law Derivation\\ of the Einstein Equations}}
\author{Robert W.\ McGwier\\
\small CTO and Staff Scientist, Cohere Technology Group, Herndon, VA}
\date{August 7, 2026 --- Version 5}

\begin{document}
\maketitle

\begin{abstract}
\noindent
We present a compact, rigorously bounded account of the claim that spacetime
curvature is induced by the entanglement structure of an underlying quantum
state, and we state precisely what has been derived, what is conjectured, and
what remains unresolved. The central verified result is that, in holographic
conformal field theories, the first law of entanglement
$\delta S_A=\delta\langle H_{\mathrm{mod},A}\rangle$, imposed on all
ball-shaped regions $A$, is \emph{equivalent} to the linearized Einstein
equations for the dual metric perturbation. An excitation---including a
quasiparticle---curves spacetime not through its intrinsic entanglement with a
partner, but through the perturbation it induces in the vacuum entanglement
across every partition. We derive the logical chain, delimit its domain of
validity, and record the negative result that pairwise (Bell-type)
entanglement between excitations carries no classical curvature signature at
this order. We further show that the framework preserves the quantum
superposition principle for geometry as a theorem rather than a postulate:
at linear order because the first law is a linear functional of the state,
and nonperturbatively because the Ryu--Takayanagi area is an operator in the
center of the region algebra, so superposed geometries are simply superposed
area eigenstates. We then extend the framework to Riemann--Cartan geometry:
the Einstein--Cartan--Sciama--Kibble equations are derived variationally and
thermodynamically, and we prove a linearized corollary that the same
entanglement first law governs the metric sector while torsion, being
algebraic in the spin density, is produced by the identical bookkeeping
wherever the state carries spin. Every material claim is labeled \pP{} (proved/derived), \pC{}
(conjectured with partial support), or \pO{} (open).
\end{abstract}

\section*{Epistemic Conventions}
Following the standing convention of this research program, every material
claim is tagged:
\begin{description}[leftmargin=2.2em,itemsep=0pt]
\item[\pP] Proved or derived within a stated framework, with published proof
or derivation in the primary literature cited at the point of use.
\item[\pC] Conjectured; supported by nontrivial consistency checks but lacking
proof.
\item[\pO] Open; neither derived nor refuted.
\end{description}
No claim is asserted without a tag or a citation. Numerical agreement is not
treated as proof. Consensus is not treated as proof.

\section{Scope and Statement of the Problem}

The question addressed is exact: \emph{in what sense, and in what regime, does
quantum entanglement produce spacetime curvature?} Three historically distinct
programs bear on it:

\begin{enumerate}[itemsep=1pt]
\item \textbf{Horizon thermodynamics.} Bekenstein \cite{Bekenstein1973} and
Hawking \cite{Hawking1975} established that a black-hole horizon of area $A$
carries entropy $S_{BH}=A/4G\hbar$ (units $c=k_B=1$). Jacobson
\cite{Jacobson1995} showed that demanding the Clausius relation
$\delta Q=T\,\delta S$, with $S\propto A$ on every local Rindler horizon and
$T$ the Unruh temperature \cite{Unruh1976}, implies the full Einstein
equations. \pP{} as a derivation \emph{given} the entropy--area postulate and a
background causal structure; the postulate itself was input, not output.
\item \textbf{Holographic entanglement entropy.} Within AdS/CFT
\cite{Maldacena1998}, Ryu and Takayanagi \cite{RyuTakayanagi2006} proposed,
and subsequent work substantiated, that the entanglement entropy of a boundary
region $A$ equals the area of the bulk minimal surface homologous to $A$:
$S_A=\mathrm{Area}(\tilde A)/4G_N$, covariantized in \cite{HRT2007} and
corrected quantum-mechanically in \cite{FLM2013,EngelhardtWall2015}. \pP{}
within AdS/CFT at leading order in $1/N$.
\item \textbf{Entropic and emergent gravity.} Verlinde \cite{Verlinde2011}
derived the Newtonian force law and the law of inertia $F=ma$ from
coarse-grained entropy gradients on holographic screens, and later argued
\cite{Verlinde2017} that in de Sitter space a volume-law entanglement
contribution modifies the gravitational law in the deep-infrared, recovering
the Milgrom acceleration scale \cite{Milgrom1983}. \pC{}: heuristic
derivations with observational motivation; no microscopic completion.
\end{enumerate}

The strongest and cleanest statement---the subject of this paper---belongs to
none of these alone. It is the \emph{first-law derivation}: linearized
gravitational dynamics is equivalent to entanglement thermodynamics
\cite{Lashkari2014,Faulkner2014}. We now derive it.

\section{Kinematics: Modular Hamiltonians and the First Law}

\begin{definition}[Reduced state and modular Hamiltonian]
Let $\rho_A=\Tr_{\bar A}\,|\Psi\rangle\langle\Psi|$ be the reduced density
matrix of a spatial region $A$. Write $\rho_A=e^{-H_A}/\Tr\,e^{-H_A}$; the
operator $H_A$ is the modular Hamiltonian. The entanglement entropy is
$S_A=-\Tr\rho_A\log\rho_A$.
\end{definition}

For a general region $H_A$ is nonlocal and unknown. Two cases are exactly
solved. \pP{}

\begin{proposition}[Bisognano--Wichmann \cite{BisognanoWichmann1976}]
For the Rindler wedge $x^1>|t|$ in the vacuum of any Wightman QFT,
$H_A=2\pi K_1$, the boost generator; equivalently
$H_A = 2\pi\!\int_{x^1>0} d^{d-1}x\; x^1\, T_{00}(x)$.
\end{proposition}

\begin{proposition}[Casini--Huerta--Myers \cite{CHM2011}]
For a ball $B$ of radius $R$ centered at $x_0$ in the vacuum of a CFT on
Minkowski space,
\begin{equation}
H_B \;=\; 2\pi\!\int_B d^{d-1}x\;\frac{R^2-|\vec x-\vec x_0|^2}{2R}\;
T_{00}(x)\,.
\label{eq:CHM}
\end{equation}
\end{proposition}

\begin{theorem}[First law of entanglement]
For any one-parameter family of states $\rho_A(\lambda)$ with
$\rho_A(0)=\rho_A^{\mathrm{vac}}$,
\begin{equation}
\delta S_A \;=\; \delta\langle H_A\rangle
\qquad\text{at first order in }\lambda .
\label{eq:firstlaw}
\end{equation}
\end{theorem}

\begin{proof}
$S(\rho_A(\lambda)) = -\Tr\rho_A\log\rho_A$. First-order variation gives
$\delta S_A = -\Tr(\delta\rho_A\log\rho_A^{\mathrm{vac}})
-\Tr(\delta\rho_A)$. The second term vanishes by trace preservation; the
first equals $\Tr(\delta\rho_A H_A)=\delta\langle H_A\rangle$ up to an additive
constant absorbed in normalization. \pP{}
\end{proof}

Equivalently: the relative entropy
$S(\rho\|\sigma)=\Tr\rho\log\rho-\Tr\rho\log\sigma$ satisfies
$S(\rho_A\|\rho_A^{\mathrm{vac}})=\Delta\langle H_A\rangle-\Delta S_A\ge 0$,
vanishing to first order \cite{BlancoCasini2013,Casini2008}. Positivity of
relative entropy is the exact quantum statement underlying both the first law
and the Bekenstein bound \cite{Casini2008}. \pP{}

\section{Dynamics: The First Law Implies Linearized Einstein Equations}

We now state the central theorem, established by Lashkari, McDermott and Van
Raamsdonk \cite{Lashkari2014} and, in full generality and both directions, by
Faulkner, Guica, Hartman, Myers and Van Raamsdonk \cite{Faulkner2014}.

\smallskip
\noindent\textbf{Setup.} Let a holographic CFT$_d$ in state
$|\Psi\rangle=|0\rangle+\lambda|\delta\Psi\rangle$ be dual to a bulk geometry
$g=g_{\mathrm{AdS}}+\lambda h$ on AdS$_{d+1}$. Two independent computations of
$\delta S_B$ exist for every ball $B$:
\begin{align}
\delta S_B^{\mathrm{grav}}
&= \frac{1}{4G_N}\,\delta\,\mathrm{Area}(\tilde B)[h]
&&\text{(Ryu--Takayanagi \cite{RyuTakayagi_note})}\,,\\
\delta\langle H_B\rangle
&= 2\pi\!\int_B d^{d-1}x\,\frac{R^2-r^2}{2R}\,
\delta\langle T_{00}\rangle
&&\text{(CHM, Eq.~\eqref{eq:CHM})}\,,
\end{align}
where $\delta\langle T_{\mu\nu}\rangle$ is read off from the asymptotics of
$h$ by the standard holographic dictionary.

\begin{theorem}[Faulkner--Guica--Hartman--Myers--Van Raamsdonk
\cite{Faulkner2014}]
The condition
\[
\delta S_B=\delta\langle H_B\rangle
\quad\text{for all balls $B$ of all radii, centers, and boost frames}
\]
holds \emph{if and only if} $h_{\mu\nu}$ satisfies the Einstein equations
linearized about AdS$_{d+1}$, sourced by $\delta\langle T_{\mu\nu}\rangle$.
\pP{}
\end{theorem}

\noindent\textbf{Structure of the proof (both directions).}
($\Leftarrow$) With $h$ on shell, the variation of the extremal area reduces,
via a bulk Noether identity of Iyer--Wald type, to a boundary term reproducing
$\delta\langle H_B\rangle$. ($\Rightarrow$) The difference
$\delta S_B-\delta\langle H_B\rangle$ can be written as the integral over the
region between $B$ and $\tilde B$ of a form $\boldsymbol{\chi}$ whose exterior
derivative is proportional to the linearized equations of motion contracted
with a Killing vector; demanding the difference vanish for the full
$(2d-1)$-parameter family of balls forces the integrand to vanish pointwise.
The constraint components follow from radius-differentiation, the dynamical
components from boosted balls. \pP{}

\begin{remark}[Interpretation]
Curvature at this order \emph{is} entanglement bookkeeping. A localized
excitation---a quasiparticle---perturbs $\rho_B$ for every ball it enters;
Eq.~\eqref{eq:firstlaw} then fixes $\delta S_B$; consistency of these
$\delta S_B$ across all partitions is possible only if an emergent metric
perturbation exists and obeys linearized Einstein dynamics sourced by the
excitation's stress tensor. The gravitational field of matter is the unique
consistent solution to vacuum-entanglement perturbation theory. \pP{} within
the stated holographic framework.
\end{remark}

\section{Beyond Linear Order}

\begin{enumerate}[itemsep=2pt]
\item \textbf{Second order.} At $O(\lambda^2)$ the first law acquires a
correction: $\Delta S_B=\Delta\langle H_B\rangle-S(\rho_B\|\rho_B^{\mathrm{vac}})$.
Positivity and monotonicity of relative entropy translate into bulk energy
conditions and constrain the second-order metric response
\cite{BlancoCasini2013,Lashkari2014}. Faulkner, Haehl, Hijano, Parrikar,
Rabideau and Van Raamsdonk \cite{Faulkner2017} showed that the CFT modular
energy at second order reproduces the \emph{nonlinear} Einstein equations to
that order, with bulk matter contributions entering through the quantum
correction $S_{\mathrm{bulk}}$ of \cite{FLM2013}. \pP{} at the perturbative
orders computed; all-orders statement \pO{}.
\item \textbf{Entanglement equilibrium.} Jacobson \cite{Jacobson2016} proposed
that the vacuum is a maximum of entanglement entropy in small causal diamonds
at fixed volume; stationarity, $\delta S_{UV}+\delta S_{IR}=0$, yields the
full nonlinear Einstein equations at each point, for arbitrary (not
necessarily AdS) backgrounds. The derivation assumes (i) an area-law UV term
with universal coefficient and (ii) a CFT-like modular response of the IR
term. \pC{}: compelling, framework-independent in aspiration, but resting on
unproven assumptions about the UV completion.
\item \textbf{Connectivity from entanglement.} Van Raamsdonk
\cite{VanRaamsdonk2010} argued, via the thermofield double, that dialing
entanglement between two CFTs to zero pinches off the Einstein--Rosen bridge:
classical connectivity of space is an entanglement effect. Maldacena and
Susskind \cite{MaldacenaSusskind2013} extended this to the conjecture ER=EPR:
any entangled pair is connected by a (generically Planckian, non-traversable)
wormhole. The thermofield-double case is \pP{} within AdS/CFT; the
extrapolation to arbitrary pairs is \pC{}.
\item \textbf{Tensor networks.} Swingle \cite{Swingle2012} observed that the
MERA tensor network representing a CFT ground state reproduces AdS geometry
with the RT prescription realized as minimal cuts, supplying a constructive
toy model of geometry-from-entanglement. \pC{} as a correspondence; exact
dictionary \pO{}.
\end{enumerate}

\section{Superposition: Why Emergent Curvature Remains Quantum}
\label{sec:superposition}

A recurrent objection to emergent-gravity programs is that a
thermodynamic or statistical origin of curvature would forfeit the quantum
superposition principle for the gravitational field. Within the present
framework the objection fails, and it fails for exact, stateable reasons.
Geometry is a functional of the quantum state; states superpose because the
Hilbert space is linear; therefore geometries superpose \emph{automatically},
without quantizing the metric as an independent field. We make this precise
at three levels.

\subsection{Linear order: the first law is a linear functional}

Both sides of Eq.~\eqref{eq:firstlaw} are linear in the perturbation
$\delta\rho_A$: $\delta\langle H_A\rangle=\Tr(\delta\rho_A H_A)$ manifestly,
and $\delta S_A=\Tr(\delta\rho_A H_A)$ by the proof of Theorem~2.1. The FGHMV
theorem therefore maps a linear space of state perturbations to a linear space
of metric perturbations: if $\delta\rho^{(1)}\mapsto h^{(1)}$ and
$\delta\rho^{(2)}\mapsto h^{(2)}$, then
$a\,\delta\rho^{(1)}+b\,\delta\rho^{(2)}\mapsto a\,h^{(1)}+b\,h^{(2)}$.
Linearized emergent gravity is a linear theory \emph{because} entanglement
perturbation theory is linear; the superposition principle for weak
gravitational fields is not an additional postulate but a corollary of the
derivation. \pP{}

Two consequences deserve emphasis. First, a quasiparticle prepared in a
superposition of two trajectories perturbs $\rho_B$ by the corresponding
superposed $\delta\rho_B$; the induced $h_{\mu\nu}$ is the superposition of
the two single-trajectory fields, including off-diagonal (coherence) data
carried by the state, not by any classical average. Second, this immediately
rules out, within the framework, the naive semiclassical prescription
$G_{\mu\nu}=8\pi G\langle T_{\mu\nu}\rangle_{\mathrm{cat}}$ applied to
macroscopic superpositions---the prescription already excluded experimentally
by Page and Geilker \cite{PageGeilker1981}. The emergent field of a cat state
is a cat state of fields, not the field of the mean mass distribution. \pP{}
within the framework; \pP{} empirically against the semiclassical alternative
\cite{PageGeilker1981}.

\subsection{Macroscopic superpositions: the area operator}

Beyond perturbation theory, superpositions of \emph{distinct classical
geometries} are governed by the quantum-error-correction formulation of
holography \cite{ADH2015,HaPPY2015}. Harlow \cite{Harlow2017} proved that in
any code subspace with complementary recovery, the entanglement entropy takes
the form
\begin{equation}
S(\rho_A)\;=\;\Tr\!\big(\rho\,\mathcal{L}_A\big)\;+\;S_{\mathrm{alg}}(\rho_A)\,,
\qquad
\mathcal{L}_A=\frac{\widehat{\mathrm{Area}}(\tilde A)}{4G_N}\,,
\label{eq:areaop}
\end{equation}
where $\mathcal{L}_A$ is an operator in the \emph{center} of the algebra of
region $A$. The Ryu--Takayanagi area is thereby promoted from a number to an
observable. A state $|\Psi\rangle=\sum_i c_i|g_i\rangle$ superposing a small
number ($\ll e^{S}$) of semiclassical geometries is an eigenvector of nothing:
it assigns amplitude $c_i$ to each area eigenvalue, and measurements of
geometric quantities obey the Born rule with respect to these amplitudes.
Distinct macroscopic geometries occupy nearly orthogonal branches
(approximate superselection by the center), which is precisely the mechanism
by which classical geometry decoheres out of the fundamentally superposed
description. \pP{} within holographic quantum error correction
\cite{ADH2015,Harlow2017}; the tensor-network models of \cite{HaPPY2015}
realize Eq.~\eqref{eq:areaop} exactly. Extension outside AdS/CFT: \pO{}.

\subsection{No collapse required, and an empirical discriminator}

Because curvature is emergent rather than fundamental, the framework requires
no gravitationally induced reduction of the state vector; it stands in direct
opposition to objective-collapse proposals in which superposed geometries are
dynamically forbidden \cite{Penrose1996}. The two classes of theory are
empirically separable: if gravity is capable of entangling two mesoscopic
masses through their mutual field---the witness protocols of Bose et al.\
\cite{Bose2017} and Marletto--Vedral \cite{MarlettoVedral2017}---then the
gravitational field sustains superposition and carries quantum information,
as the entanglement-first framework requires. A confirmed
gravitationally-mediated-entanglement signal would be consistent with, though
not uniquely probative of, geometry-from-entanglement; a confirmed null
result at sufficient sensitivity would falsify it. \pC{} pending experiment;
the logical structure of the discriminator is \pP{}.

\section{GR plus Torsion: The Einstein--Cartan Extension, Derived}
\label{sec:torsion}

We now extend the program to Riemann--Cartan geometry, in which the
connection carries torsion. The extension is executed in three steps: the
classical variational derivation \pP{}, the thermodynamic (Clausius)
derivation \pP{} within stated assumptions, and an entanglement first-law
corollary derived here, with its unproven remainder tagged honestly.

\subsection{Variational derivation of Einstein--Cartan--Sciama--Kibble
theory}

Work in first-order variables: the coframe $e^a=e^a{}_\mu dx^\mu$ and an
independent Lorentz connection $\omega^{ab}=\omega^{ab}{}_\mu dx^\mu$, with
\begin{equation}
T^a \equiv de^a+\omega^a{}_b\wedge e^b\,,\qquad
R^{ab}\equiv d\omega^{ab}+\omega^a{}_c\wedge\omega^{cb}\,,
\end{equation}
the torsion and curvature two-forms. Take the Palatini action in $d=4$
($\kappa=8\pi G$),
\begin{equation}
S \;=\; \frac{1}{4\kappa}\int \epsilon_{abcd}\,e^a\wedge e^b\wedge R^{cd}
\;+\; S_m[e,\omega,\psi]\,.
\label{eq:ECaction}
\end{equation}
Independent variation gives two field equations \cite{Kibble1961,Sciama1964,
HehlRMP1976,Trautman2006}:
\begin{align}
\delta e:&\quad
\tfrac{1}{2}\,\epsilon_{abcd}\,e^b\wedge R^{cd} \;=\; \kappa\,\tau_a\,,
\label{eq:ECeinstein}\\[2pt]
\delta\omega:&\quad
\epsilon_{abcd}\,e^a\wedge T^b \;=\; \kappa\,\mathcal{S}_{cd}\,,
\label{eq:cartan}
\end{align}
where $\tau_a$ is the canonical energy--momentum three-form (generally
asymmetric) and $\mathcal{S}_{cd}=2\,\delta S_m/\delta\omega^{cd}$ the
canonical spin-current three-form. In tensor components, with the
conventions of \cite{HehlRMP1976}, Eq.~\eqref{eq:cartan} reads
\begin{equation}
T^\lambda{}_{\mu\nu}
+\delta^\lambda_\mu\,T^\sigma{}_{\nu\sigma}
-\delta^\lambda_\nu\,T^\sigma{}_{\mu\sigma}
\;=\;\kappa\, s^\lambda{}_{\mu\nu}\,,
\label{eq:cartantensor}
\end{equation}
the \emph{Cartan equation}: torsion is determined \emph{algebraically},
pointwise, by the spin density. It contains no derivatives of $T$; torsion
does not propagate, and in vacuum ($s=0$) it vanishes identically, whereupon
$\omega$ reduces to the Levi-Civita connection and
Eq.~\eqref{eq:ECeinstein} reduces to the Einstein equations. \pP{}

\smallskip
\noindent\textbf{Elimination of torsion.} Decompose
$\omega=\tilde\omega(e)+K$ with $\tilde\omega$ Levi-Civita and $K$ the
contortion one-form, fixed linearly in $s$ by
Eq.~\eqref{eq:cartantensor}. Substituting into
Eqs.~\eqref{eq:ECaction}--\eqref{eq:ECeinstein} yields ordinary Riemannian
GR sourced by the symmetric Belinfante--Rosenfeld stress tensor, plus a
contact term quadratic in the spin density. For a minimally coupled Dirac
field this is the Hehl--Datta interaction \cite{HehlDatta1971},
\begin{equation}
\mathcal{L}_{\mathrm{eff}}
=\mathcal{L}_{\mathrm{GR}}+\mathcal{L}_{\mathrm{Dirac}}
-\frac{3\kappa}{16}\,
\big(\bar\psi\gamma^5\gamma^\mu\psi\big)
\big(\bar\psi\gamma^5\gamma_\mu\psi\big)\,,
\label{eq:hehldatta}
\end{equation}
an axial--axial four-fermion coupling suppressed by $G$, negligible except
where $\kappa\,s^2$ is comparable to the energy density---far beyond any
astrophysically realized regime, though decisive near classical
singularities, where the induced repulsion can avert collapse
\cite{Trautman1973}. \pP{} for the derivation;
singularity aversion in realistic cosmology \pC{} \cite{Poplawski2010}.

\subsection{Thermodynamic derivation with torsion}

The Jacobson construction of Section~1 extends to Riemann--Cartan
spacetime. Dey, Liberati and Pranzetti \cite{DLP2017} imposed the Clausius
relation $\delta Q=T\,\delta S$ on local causal horizons of a
Riemann--Cartan geometry, with the entropy given by the Noether charge of
the first-order theory, and recovered precisely the Einstein--Cartan
equations \eqref{eq:ECeinstein}--\eqref{eq:cartan}: the algebraic Cartan
sector rides along with the entropy balance rather than obstructing it. The
treatment of the contorsion contribution to the entropy was subsequently
sharpened by De Lorenzo, De Paoli and Speziale \cite{DDS2018}. That the
horizon entropy remains $\mathrm{Area}/4G$ in first-order variables---i.e.,
that torsion does not renormalize the entropy density for the
Einstein--Cartan action---is established by the Lorentz-diffeomorphism
Noether-charge analysis of Jacobson and Mohd \cite{JacobsonMohd2015}. \pP{}
within the stated assumptions (local equilibrium, Noether-charge entropy);
these are the same assumptions as the torsion-free case plus regularity of
the internal Lorentz frame on the horizon.

\subsection{Entanglement first law with torsion: a corollary and its
remainder}

We now state what the first-law mechanism of Section~3 implies for
Einstein--Cartan gravity. The derivation below is presented here; its
elementary steps are \pP{}, and the unproven remainder is tagged.

\begin{proposition}[Linearized decomposition]
Linearize Einstein--Cartan theory about a maximally symmetric vacuum
(AdS$_{d+1}$ or Minkowski). Then: (i) the background torsion vanishes,
because the vacuum spin current vanishes by Lorentz invariance and
Eq.~\eqref{eq:cartantensor} is algebraic; (ii) at first order the system
decomposes exactly into a metric sector, in which $h_{\mu\nu}$ obeys the
linearized \emph{Einstein} equations sourced by the Belinfante--Rosenfeld
$\delta\langle T_{\mu\nu}\rangle$, and an algebraic torsion sector, in which
the contortion perturbation is fixed pointwise by $\delta\langle
s^\lambda{}_{\mu\nu}\rangle$ with no independent dynamics; (iii) the
Hehl--Datta contact term \eqref{eq:hehldatta}, being quadratic in $s$,
contributes nothing at this order. \pP{} (each step is an inspection of
Eqs.~\eqref{eq:ECeinstein}--\eqref{eq:hehldatta}).
\end{proposition}

\begin{proposition}[First-law equivalence, metric sector]
Under the hypotheses of the FGHMV theorem (Section~3), and using that the
CFT stress tensor entering the Casini--Huerta--Myers modular Hamiltonian
\eqref{eq:CHM} is the symmetric (Belinfante) one, the first law
$\delta S_B=\delta\langle H_B\rangle$ on all balls is equivalent to the
linearized Einstein--Cartan equations in the metric sector. The torsion
sector adds no independent first-law content at this order: carrying no
propagating degrees of freedom, it contributes to $\delta S_B$ only through
the matter state perturbation $\delta\rho_B$ already accounted for. \pP{}
as a corollary of \cite{Faulkner2014} plus the decomposition above.
\end{proposition}

\begin{remark}[The honest remainder]
What is \emph{not} established: (i) a boundary first law in the presence of
an explicit source coupled to the spin current $s^\lambda{}_{\mu\nu}$,
which would deform the modular Hamiltonian away from the CHM form---the
holographic dictionary for such torsionful boundary conditions has been
developed in special settings \cite{GallegosGursoy2020} but no
FGHMV-grade theorem exists; \pO{} (ii) any statement at second order, where
the Hehl--Datta term activates and the bulk entanglement correction
$S_{\mathrm{bulk}}$ \cite{FLM2013} must be evaluated on a torsionful
background; \pO{} (iii) propagating-torsion extensions (general Poincar\'e
gauge theory), which change the count of degrees of freedom and void the
decomposition of the Proposition; \pO{}. Within
Einstein--Cartan proper---torsion algebraic, sourced by spin---the verdict
is clean: \emph{entanglement produces curvature exactly as before, and the
same entanglement bookkeeping, read through the Cartan equation, produces
torsion wherever the state carries spin density.} \pP{} at linear order.
\end{remark}

\section{Negative Results: What Entanglement Does \emph{Not} Do}

Precision requires stating what the theorems do not license.

\begin{enumerate}[itemsep=2pt]
\item \textbf{Pair entanglement does not gravitate.} The curvature sourced by
a quasiparticle in Section~3 is fixed entirely by
$\delta\langle T_{\mu\nu}\rangle$. Two excitations prepared in a Bell state
versus a product state with the same stress-tensor profile produce the
\emph{same} linearized metric. Mutual entanglement between excitations is a
microstate correlation, not a source term. \pP{} at linear order, by
inspection of the theorem's content. This is consistent with
Section~\ref{sec:superposition}: the metric responds to $\delta\rho_B$ ball
by ball; a single-particle superposition of trajectories \emph{changes} every
$\delta\rho_B$ and hence the field, whereas entangling two particles without
altering their single-region reduced states does not. Whether ER=EPR endows such pairs with a
Planckian geometric dual is \pC{}; even if so, it carries no classical
curvature signature accessible to external measurement.
\item \textbf{Coarse-grained entropic forces are a distinct claim.}
Verlinde's 2010 derivation \cite{Verlinde2011} requires $N\gg1$ bits on a
screen at finite temperature; it is undefined for a single quasiparticle and
must not be conflated with the exact first-law mechanism. \pP{} (a statement
about scope).
\item \textbf{Analog gravity is disanalogous.} In condensed-matter analogs
(Unruh's sonic metric \cite{Unruh1981}; review \cite{BLV2005}), quasiparticles
propagate on an effective curved metric determined by the background
condensate flow. The analog Einstein equations are \emph{not} satisfied, and
quasiparticle entanglement plays no dynamical geometric role. The
entanglement--curvature mechanism is specific to vacua whose leading
entanglement is area-law with the gravitational coefficient $1/4G_N$. \pP{}
\item \textbf{De Sitter is unresolved.} All \pP{} results above are anchored
in AdS/CFT. Our universe has positive cosmological constant; the relevant
entropy includes horizon and possibly volume-law contributions
\cite{Verlinde2017}, quantum extremal surface technology
\cite{EngelhardtWall2015} lacks a settled de Sitter formulation, and
Susskind's dictum that ``entanglement is not enough'' points to computational
complexity as an additional ingredient. \pO{}
\end{enumerate}

\section{Certified Conclusions}

\begin{enumerate}[itemsep=1pt]
\item In holographic CFTs, linearized Einstein gravity---curvature sourced by
stress-energy---is mathematically equivalent to the first law of entanglement
imposed on all ball-shaped regions. \pP{} \cite{Lashkari2014,Faulkner2014}
\item The mechanism by which an excitation curves spacetime is its
perturbation of the \emph{vacuum} entanglement structure, quantified by
$\delta S_B=\delta\langle H_B\rangle$, not its intrinsic entanglement with
other excitations. \pP{}
\item Nonlinear corrections through second order follow from relative-entropy
identities and bulk modular flow. \pP{} at computed orders
\cite{Faulkner2017}; complete nonperturbative statement \pO{}.
\item The superposition principle for gravity is derived, not assumed:
linearity of the first law makes weak-field superposition a corollary
\pP{}, and the area operator of holographic error correction
\cite{Harlow2017} makes superpositions of macroscopically distinct
curvatures well-defined states with Born-rule statistics \pP{} (within
AdS/CFT; beyond it \pO{}). No gravitational collapse mechanism is required,
and gravitationally mediated entanglement experiments
\cite{Bose2017,MarlettoVedral2017} furnish an empirical discriminator
against collapse alternatives \cite{Penrose1996}. \pC{} pending data.
\item GR plus torsion is accommodated, not merely conjectured:
Einstein--Cartan theory follows variationally \pP{}
\cite{Kibble1961,Sciama1964,HehlRMP1976}, thermodynamically \pP{}
\cite{DLP2017,DDS2018,JacobsonMohd2015}, and, at linear order about the
vacuum, from the entanglement first law in the metric sector with torsion
fixed algebraically by the spin density of the state \pP{} (this work,
Section~\ref{sec:torsion}). Spin-current-deformed modular Hamiltonians,
second-order torsion effects, and propagating-torsion generalizations are
\pO{}.
\item Framework-independent versions (entanglement equilibrium, ER=EPR,
emergent de Sitter gravity) are live, partially supported, and unproven.
\pC{}/\pO{} as itemized above.
\end{enumerate}

No claim above rests on consensus, numerology, or agreement among secondary
sources. Each \pP{} tag corresponds to a published derivation verified against
the primary reference. Where the chain of derivation ends, the tag says so.

\appendix

\section{Domain of Validity: AdS versus de Sitter}
\label{app:adsds}

This appendix supplies the technical basis for the exclusive use of
anti-de Sitter holography in the \pP{}-grade results of the main text,
despite the observed positive cosmological constant, together with the
precise argument for why those results nevertheless govern local physics
in our universe.

\subsection{Why the proofs live in AdS}

The \pP{}-grade chain---CHM modular Hamiltonians \cite{CHM2011}, the first
law, and the FGHMV equivalence \cite{Faulkner2014,Lashkari2014}---requires
three ingredients that AdS supplies and de Sitter does not:

\begin{enumerate}[itemsep=2pt]
\item \textbf{A boundary anchoring the microscopic theory.} AdS is
conformally a box: a timelike boundary at spatial infinity carries a
complete, unitary CFT, and the Maldacena dictionary \cite{Maldacena1998}
defines every bulk quantity as a boundary observable. De Sitter has no
such boundary; its candidate asymptotics are spacelike, and the proposed
dS/CFT correspondence \cite{Strominger2001} pairs the bulk with a
Euclidean, generically non-unitary boundary theory whose status remains
\pO{} after two decades.
\item \textbf{A global timelike Killing vector} (in the relevant patches),
guaranteeing a preferred vacuum, positive energy, and the applicability of
the Bisognano--Wichmann machinery. De Sitter possesses only
observer-dependent static patches with horizons; the finite
Gibbons--Hawking horizon entropy \cite{GibbonsHawking1977} suggests a
finite-dimensional Hilbert space per observer, structurally unlike the
boundary CFT, and no consensus microscopic description exists. \pO{}
\item \textbf{Theorem-grade entanglement technology.} Ryu--Takayanagi and
its quantum extensions
\cite{RyuTakayanagi2006,FLM2013,EngelhardtWall2015} are proved for AdS
boundary conditions; no de Sitter analogue has reached that grade. \pO{}
\end{enumerate}

The methodological rule of the main text follows: certify in AdS,
extrapolate with labels.

\subsection{Why the certified mechanism nevertheless governs our space}

\begin{proposition}[Link 1: flat-space inputs]
The quantum-informational inputs of the mechanism are theorems of QFT on
Minkowski space, independent of holography: the CHM ball modular
Hamiltonian \cite{CHM2011}, the first law
$\delta S=\delta\langle H\rangle$ (Theorem 2.1), positivity and
monotonicity of relative entropy \cite{Casini2008,BlancoCasini2013}, and
the Markov property of the vacuum on null planes
\cite{CasiniTesteTorroba2017}. The vacuum of the Standard-Model fields in
our universe therefore carries the entanglement structure the mechanism
runs on. \pP{}
\end{proposition}

\begin{proposition}[Link 2: locality of the derivation]
The thermodynamic route \cite{Jacobson1995,DLP2017} uses only local
Rindler horizons---available through every point of every spacetime---and
yields the Einstein(--Cartan) equations pointwise. AdS/CFT is what
upgrades the underlying entropy--area input from assumption to theorem;
the assumptions so validated are local and are satisfied in our universe
wherever curvature is small compared to microscopic scales. \pP{} for the
structure; the transfer to general backgrounds is the content of
entanglement equilibrium \cite{Jacobson2016} and is \pC{}.
\end{proposition}

\begin{proposition}[Link 3: infrared decoupling of $\Lambda$]
The observed cosmological constant is
$\Lambda\simeq1.1\times10^{-52}\,\mathrm{m}^{-2}$. The dimensionless
measure of its local relevance at scale $L$ is $\Lambda L^2$:
$\sim10^{-52}$ (laboratory), $\sim10^{-30}$ (1 AU), $\sim10^{-11}$
(10 kpc), $\mathcal{O}(1)$ only at the Hubble radius. Below the Hubble
scale the AdS/dS/flat distinction is invisible to any local experiment to
the stated precision; the certified linearized regime is the
$\Lambda L^2\ll1$ regime. \pP{} (arithmetic on the measured $\Lambda$).
\end{proposition}

\begin{proposition}[Link 4: the difference bites where the open problems
are filed]
The sign of $\Lambda$ becomes structural only near the horizon scale,
where de Sitter has what AdS lacks: an observer horizon with entropy
$A/4G$ \cite{GibbonsHawking1977} and, in Verlinde's proposal, a competing
volume-law entanglement contribution \cite{Verlinde2017}. Any resulting
modification appears at accelerations of order $cH_0$; numerically,
Milgrom's $a_0\simeq1.2\times10^{-10}\,\mathrm{m\,s^{-2}}$ satisfies
$a_0\approx cH_0/6$ \cite{Milgrom1983}, the regime of the observed
rotation-curve anomalies. The coincidence is \pP{} as arithmetic; its
explanation by horizon-scale entanglement is \pC{}; the underlying dS
microtheory is \pO{}.
\end{proposition}

\section{The Status of Wick Rotation}
\label{app:wick}

Euclidean continuation enters the program at three load-bearing points:
the path-integral representation of reduced density matrices, the
identification of modular flow, and the anomaly coefficient in the Gauge
Lemma of the companion note \cite{McGwier2026b}. Each is covered by a
theorem, and the historically fragile uses of Wick rotation are absent
from every \pP{} claim.

\subsection{Reconstruction and instantiation}

\begin{theorem}[Osterwalder--Schrader \cite{OS1973,OS1975}]
Euclidean correlation functions satisfying reflection positivity,
Euclidean covariance, symmetry, and cluster/regularity conditions
determine, uniquely, a Lorentzian QFT satisfying the Wightman axioms:
unitarity, locality, positive energy, unique vacuum.
\end{theorem}

The boundary theories used at \pP{} grade satisfy the hypotheses; every
Euclidean computation of modular data converts to real-time physics by
theorem, not analogy. \pP{} The rotation is then \emph{instantiated} once,
at the foundational level, by Bisognano--Wichmann
\cite{BisognanoWichmann1976} (Proposition 2.2): modular flow of the wedge
\emph{is} the boost, with the CHM ball result its conformal image.
Tomita--Takesaki modular theory---the source of the first law and
relative-entropy positivity---is formulated directly in the Lorentzian
operator algebra \cite{Haag1996} and performs no rotation at all. The
physical countersignature is the Unruh temperature $T=\hbar a/2\pi ck_B$
\cite{Unruh1976}: thermality under modular flow is the KMS condition, the
operational fingerprint of valid Euclidean periodicity. \pP{}

\subsection{The fragile uses, and their absence from \pP{} claims}

\begin{enumerate}[itemsep=2pt]
\item \textbf{Replica continuation $n\to1$.} The Lewkowycz--Maldacena
derivation of RT \cite{LM2013} continues in the number of replica sheets,
requiring assumptions beyond Osterwalder--Schrader. The central
linear-order result of the main text does not: the first law is an
operator identity (Theorem 2.1) needing no replicas, and FGHMV
\cite{Faulkner2014} uses RT only as an independently established input.
Where replica technology is genuinely load-bearing
\cite{FLM2013,EngelhardtWall2015}, the corresponding claims carry their
perturbative qualifications explicitly. \pP{} for the accounting.
\item \textbf{The Euclidean gravitational path integral.} The Euclidean
Einstein--Hilbert action is unbounded below (the conformal-factor problem
\cite{GHP1978}); integrals over Euclidean metrics as fundamental variables
are formal. The program never performs one: gravity is emergent, and every
gravitational statement is a statement about boundary entanglement. The
fragility is bypassed structurally, not resolved. \pP{} (that it is
bypassed); Euclidean quantum gravity itself remains \pO{} and is not
relied upon.
\item \textbf{De Sitter Euclidean vacua.} Continuation in de Sitter
(sphere partition functions, $\alpha$-vacua) is genuinely subtle for lack
of a global timelike Killing vector. This fragility lands entirely inside
the region already tagged \pO{} in Appendix~\ref{app:adsds}: the two
frontiers coincide---a consistency check on the labeling, not an
additional exposure. \pP{}
\end{enumerate}

\subsection{Fermion conventions}

Euclidean fermions require hermiticity conventions for $\gamma^5$ and the
measure. The single certified claim this touches---the anomaly caveat in
the Gauge Lemma of \cite{McGwier2026b}---is convention-independent: the
anomaly coefficient is fixed by the Atiyah--Singer index theorem, and
Fujikawa's Euclidean evaluation \cite{Fujikawa1979} reproduces the
Lorentzian Ward identity exactly. Convention, not physics. \pP{}

\section*{Acknowledgments}
The author acknowledges the use of Claude (Anthropic, claude.ai) for
assistance in drafting, literature synthesis, and \LaTeX{} preparation of
this manuscript. All technical claims, epistemic classifications, and the
final text were reviewed and verified by the author, who bears sole
responsibility for the content.

\section*{Version History}
\begin{description}[leftmargin=2.2em,itemsep=0pt]
\item[v1] August 6, 2026. Initial release: first-law derivation, nonlinear
extensions, negative results.
\item[v2] August 6, 2026. Added Section~\ref{sec:superposition} on the
superposition principle for emergent geometry (linearity of the first law;
area operator of holographic quantum error correction; empirical
discriminators), corresponding conclusions, seven references, and
acknowledgments.
\item[v3] August 6, 2026. Added Section~\ref{sec:torsion}: variational and
thermodynamic derivations of Einstein--Cartan--Sciama--Kibble theory, a
linearized first-law corollary for the torsionful case with its open
remainder itemized, corresponding conclusions, and ten references.
\item[v4] August 6, 2026. Added Appendices~\ref{app:adsds}
and~\ref{app:wick}: the AdS-versus-de-Sitter domain-of-validity argument
in four propositions, and the theorem-grade status of Wick rotation with
its fragile uses shown absent from all \pP{} claims; nine references.
\item[v5] August 7, 2026. Program bibliography completed: all six works of
the series listed with DOIs where assigned.
\end{description}

\begin{thebibliography}{99}\small

\bibitem{Bekenstein1973}
J.~D.~Bekenstein, ``Black holes and entropy,''
Phys.\ Rev.\ D \textbf{7}, 2333 (1973).

\bibitem{Hawking1975}
S.~W.~Hawking, ``Particle creation by black holes,''
Commun.\ Math.\ Phys.\ \textbf{43}, 199 (1975).

\bibitem{Unruh1976}
W.~G.~Unruh, ``Notes on black hole evaporation,''
Phys.\ Rev.\ D \textbf{14}, 870 (1976).

\bibitem{BisognanoWichmann1976}
J.~J.~Bisognano and E.~H.~Wichmann, ``On the duality condition for quantum
fields,'' J.\ Math.\ Phys.\ \textbf{17}, 303 (1976).

\bibitem{Milgrom1983}
M.~Milgrom, ``A modification of the Newtonian dynamics as a possible
alternative to the hidden mass hypothesis,''
Astrophys.\ J.\ \textbf{270}, 365 (1983).

\bibitem{Unruh1981}
W.~G.~Unruh, ``Experimental black-hole evaporation?,''
Phys.\ Rev.\ Lett.\ \textbf{46}, 1351 (1981).

\bibitem{Jacobson1995}
T.~Jacobson, ``Thermodynamics of spacetime: The Einstein equation of state,''
Phys.\ Rev.\ Lett.\ \textbf{75}, 1260 (1995), arXiv:gr-qc/9504004.

\bibitem{Maldacena1998}
J.~M.~Maldacena, ``The large-$N$ limit of superconformal field theories and
supergravity,'' Adv.\ Theor.\ Math.\ Phys.\ \textbf{2}, 231 (1998),
arXiv:hep-th/9711200.

\bibitem{RyuTakayanagi2006}
S.~Ryu and T.~Takayanagi, ``Holographic derivation of entanglement entropy
from AdS/CFT,'' Phys.\ Rev.\ Lett.\ \textbf{96}, 181602 (2006),
arXiv:hep-th/0603001.

\bibitem{HRT2007}
V.~E.~Hubeny, M.~Rangamani and T.~Takayanagi, ``A covariant holographic
entanglement entropy proposal,'' JHEP \textbf{07}, 062 (2007),
arXiv:0705.0016.

\bibitem{Casini2008}
H.~Casini, ``Relative entropy and the Bekenstein bound,''
Class.\ Quantum Grav.\ \textbf{25}, 205021 (2008), arXiv:0804.2182.

\bibitem{VanRaamsdonk2010}
M.~Van Raamsdonk, ``Building up spacetime with quantum entanglement,''
Gen.\ Rel.\ Grav.\ \textbf{42}, 2323 (2010), arXiv:1005.3035.

\bibitem{Verlinde2011}
E.~Verlinde, ``On the origin of gravity and the laws of Newton,''
JHEP \textbf{04}, 029 (2011), arXiv:1001.0785.

\bibitem{CHM2011}
H.~Casini, M.~Huerta and R.~C.~Myers, ``Towards a derivation of holographic
entanglement entropy,'' JHEP \textbf{05}, 036 (2011), arXiv:1102.0440.

\bibitem{Swingle2012}
B.~Swingle, ``Entanglement renormalization and holography,''
Phys.\ Rev.\ D \textbf{86}, 065007 (2012), arXiv:0905.1317.

\bibitem{MaldacenaSusskind2013}
J.~Maldacena and L.~Susskind, ``Cool horizons for entangled black holes,''
Fortsch.\ Phys.\ \textbf{61}, 781 (2013), arXiv:1306.0533.

\bibitem{BlancoCasini2013}
D.~D.~Blanco, H.~Casini, L.-Y.~Hung and R.~C.~Myers, ``Relative entropy and
holography,'' JHEP \textbf{08}, 060 (2013), arXiv:1305.3182.

\bibitem{FLM2013}
T.~Faulkner, A.~Lewkowycz and J.~Maldacena, ``Quantum corrections to
holographic entanglement entropy,'' JHEP \textbf{11}, 074 (2013),
arXiv:1307.2892.

\bibitem{Lashkari2014}
N.~Lashkari, M.~B.~McDermott and M.~Van Raamsdonk, ``Gravitational dynamics
from entanglement `thermodynamics','' JHEP \textbf{04}, 195 (2014),
arXiv:1308.3716.

\bibitem{Faulkner2014}
T.~Faulkner, M.~Guica, T.~Hartman, R.~C.~Myers and M.~Van Raamsdonk,
``Gravitation from entanglement in holographic CFTs,''
JHEP \textbf{03}, 051 (2014), arXiv:1312.7856.

\bibitem{EngelhardtWall2015}
N.~Engelhardt and A.~C.~Wall, ``Quantum extremal surfaces: holographic
entanglement entropy beyond the classical regime,''
JHEP \textbf{01}, 073 (2015), arXiv:1408.3203.

\bibitem{Jacobson2016}
T.~Jacobson, ``Entanglement equilibrium and the Einstein equation,''
Phys.\ Rev.\ Lett.\ \textbf{116}, 201101 (2016), arXiv:1505.04753.

\bibitem{Faulkner2017}
T.~Faulkner, F.~M.~Haehl, E.~Hijano, O.~Parrikar, C.~Rabideau and
M.~Van Raamsdonk, ``Nonlinear gravity from entanglement in conformal field
theories,'' JHEP \textbf{08}, 057 (2017), arXiv:1705.03026.

\bibitem{Verlinde2017}
E.~Verlinde, ``Emergent gravity and the dark universe,''
SciPost Phys.\ \textbf{2}, 016 (2017), arXiv:1611.02269.

\bibitem{BLV2005}
C.~Barcel\'o, S.~Liberati and M.~Visser, ``Analogue gravity,''
Living Rev.\ Rel.\ \textbf{8}, 12 (2005), arXiv:gr-qc/0505065.

\bibitem{PageGeilker1981}
D.~N.~Page and C.~D.~Geilker, ``Indirect evidence for quantum gravity,''
Phys.\ Rev.\ Lett.\ \textbf{47}, 979 (1981).

\bibitem{Penrose1996}
R.~Penrose, ``On gravity's role in quantum state reduction,''
Gen.\ Rel.\ Grav.\ \textbf{28}, 581 (1996).

\bibitem{ADH2015}
A.~Almheiri, X.~Dong and D.~Harlow, ``Bulk locality and quantum error
correction in AdS/CFT,'' JHEP \textbf{04}, 163 (2015), arXiv:1411.7041.

\bibitem{HaPPY2015}
F.~Pastawski, B.~Yoshida, D.~Harlow and J.~Preskill, ``Holographic quantum
error-correcting codes: toy models for the bulk/boundary correspondence,''
JHEP \textbf{06}, 149 (2015), arXiv:1503.06237.

\bibitem{Harlow2017}
D.~Harlow, ``The Ryu--Takayanagi formula from quantum error correction,''
Commun.\ Math.\ Phys.\ \textbf{354}, 865 (2017), arXiv:1607.03901.

\bibitem{Bose2017}
S.~Bose, A.~Mazumdar, G.~W.~Morley, H.~Ulbricht, M.~Toro\v{s},
M.~Paternostro, A.~A.~Geraci, P.~F.~Barker, M.~S.~Kim and G.~Milburn,
``Spin entanglement witness for quantum gravity,''
Phys.\ Rev.\ Lett.\ \textbf{119}, 240401 (2017), arXiv:1707.06050.

\bibitem{MarlettoVedral2017}
C.~Marletto and V.~Vedral, ``Gravitationally induced entanglement between two
massive particles is sufficient evidence of quantum effects in gravity,''
Phys.\ Rev.\ Lett.\ \textbf{119}, 240402 (2017), arXiv:1707.06036.

\bibitem{Kibble1961}
T.~W.~B.~Kibble, ``Lorentz invariance and the gravitational field,''
J.\ Math.\ Phys.\ \textbf{2}, 212 (1961).

\bibitem{Sciama1964}
D.~W.~Sciama, ``The physical structure of general relativity,''
Rev.\ Mod.\ Phys.\ \textbf{36}, 463 (1964); erratum \textbf{36}, 1103
(1964).

\bibitem{HehlDatta1971}
F.~W.~Hehl and B.~K.~Datta, ``Nonlinear spinor equation and asymmetric
connection in general relativity,''
J.\ Math.\ Phys.\ \textbf{12}, 1334 (1971).

\bibitem{Trautman1973}
A.~Trautman, ``Spin and torsion may avert gravitational singularities,''
Nature Phys.\ Sci.\ \textbf{242}, 7 (1973).

\bibitem{HehlRMP1976}
F.~W.~Hehl, P.~von der Heyde, G.~D.~Kerlick and J.~M.~Nester, ``General
relativity with spin and torsion: foundations and prospects,''
Rev.\ Mod.\ Phys.\ \textbf{48}, 393 (1976).

\bibitem{Trautman2006}
A.~Trautman, ``Einstein--Cartan theory,'' in \emph{Encyclopedia of
Mathematical Physics}, eds.\ J.-P.~Fran\c{c}oise et al.\ (Elsevier, 2006),
arXiv:gr-qc/0606062.

\bibitem{Poplawski2010}
N.~J.~Pop\l awski, ``Cosmology with torsion: an alternative to cosmic
inflation,'' Phys.\ Lett.\ B \textbf{694}, 181 (2010), arXiv:1007.0587.

\bibitem{JacobsonMohd2015}
T.~Jacobson and A.~Mohd, ``Black hole entropy and Lorentz-diffeomorphism
Noether charge,'' Phys.\ Rev.\ D \textbf{92}, 124010 (2015),
arXiv:1507.01054.

\bibitem{DLP2017}
R.~Dey, S.~Liberati and D.~Pranzetti, ``Spacetime thermodynamics in the
presence of torsion,'' Phys.\ Rev.\ D \textbf{96}, 124032 (2017),
arXiv:1709.04031.

\bibitem{DDS2018}
T.~De Lorenzo, E.~De Paoli and S.~Speziale, ``Spacetime thermodynamics with
contorsion,'' Phys.\ Rev.\ D \textbf{98}, 064053 (2018), arXiv:1807.02041.

\bibitem{GallegosGursoy2020}
A.~D.~Gallegos and U.~G\"ursoy, ``Holographic spin liquids and Lovelock
Chern--Simons gravity,'' JHEP \textbf{11}, 151 (2020), arXiv:2004.05148.

\bibitem{OS1973}
K.~Osterwalder and R.~Schrader, ``Axioms for Euclidean Green's
functions,'' Commun.\ Math.\ Phys.\ \textbf{31}, 83 (1973); II,
Commun.\ Math.\ Phys.\ \textbf{42}, 281 (1975).
\bibitem{OS1975}
See \cite{OS1973}, part II.

\bibitem{GibbonsHawking1977}
G.~W.~Gibbons and S.~W.~Hawking, ``Cosmological event horizons,
thermodynamics, and particle creation,''
Phys.\ Rev.\ D \textbf{15}, 2738 (1977).

\bibitem{GHP1978}
G.~W.~Gibbons, S.~W.~Hawking and M.~J.~Perry, ``Path integrals and the
indefiniteness of the gravitational action,''
Nucl.\ Phys.\ B \textbf{138}, 141 (1978).

\bibitem{Fujikawa1979}
K.~Fujikawa, ``Path-integral measure for gauge-invariant fermion
theories,'' Phys.\ Rev.\ Lett.\ \textbf{42}, 1195 (1979).

\bibitem{Haag1996}
R.~Haag, \emph{Local Quantum Physics: Fields, Particles, Algebras},
2nd ed.\ (Springer, Berlin, 1996).

\bibitem{Strominger2001}
A.~Strominger, ``The dS/CFT correspondence,''
JHEP \textbf{10}, 034 (2001), arXiv:hep-th/0106113.

\bibitem{LM2013}
A.~Lewkowycz and J.~Maldacena, ``Generalized gravitational entropy,''
JHEP \textbf{08}, 090 (2013), arXiv:1304.4926.

\bibitem{CasiniTesteTorroba2017}
H.~Casini, E.~Test\'e and G.~Torroba, ``Modular Hamiltonians on the null
plane and the Markov property of the vacuum state,''
J.\ Phys.\ A \textbf{50}, 364001 (2017), arXiv:1703.10656.

\bibitem{McGwier2026b}
R.~W.~McGwier, ``Spin-current sources and ball modular Hamiltonians:
resolving the torsionful first law,''
Zenodo (2026), DOI:
\href{https://doi.org/10.5281/zenodo.21824934}{10.5281/zenodo.21824934}.

\bibitem{McGwier2026c}
R.~W.~McGwier, ``The New Standard Model for unified physics,''
Zenodo (2026), DOI:
\href{https://doi.org/10.5281/zenodo.21829536}{10.5281/zenodo.21829536}.

\bibitem{McGwier2026d}
R.~W.~McGwier, ``Condition NL for finite balls: a spectral reduction, the unconditional metric sector, and a universal kinematic residue,''
Zenodo (2026), DOI:
\href{https://doi.org/10.5281/zenodo.21836572}{10.5281/zenodo.21836572}.

\bibitem{McGwier2026e}
R.~W.~McGwier, ``The Hehl--Datta coefficient from relative-entropy positivity,''
Zenodo (2026), DOI:
\href{https://doi.org/10.5281/zenodo.21836757}{10.5281/zenodo.21836757}.

\bibitem{McGwier2026f}
R.~W.~McGwier, ``Emergent gravity and a new Lagrangian: a theory of everything we know,''
Zenodo (2026), DOI to be assigned.

\bibitem{RyuTakayagi_note}
Reference \cite{RyuTakayanagi2006}; covariant form \cite{HRT2007}.

\end{thebibliography}

\end{document}
